\section{گام ششم: تحلیل پیچیدگی محاسباتی و نتیجه‌گیری نهایی}

در این گام، به عنوان آخرین بخش از پروژه، مدل‌های توسعه‌یافته از منظر پیچیدگی محاسباتی، زمان آموزش و تعادل بین دقت و هزینه‌ها مورد ارزیابی قرار می‌گیرند.

\subsection{محیط سخت‌افزاری و نرم‌افزاری}
با توجه به نیاز به پردازش‌های سنگین ماتریسی و به منظور تسریع فرآیند، آموزش روی بستر ابری \lr{Google Colab} اجرا گردید. مشخصات این محیط محاسباتی به شرح زیر است:
\begin{itemize}
    \item \textbf{پردازنده‌ی گرافیکی \lr{(Device)}:} \lr{NVIDIA T4 Tensor Core GPU}
    \item \textbf{حافظه‌ی گرافیکی \lr{(GPU Memory)}:} \lr{15 GB}
    \item \textbf{حافظه‌ی اصلی \lr{(RAM)}:} \lr{12.7 GB}
    \item \textbf{نسخه‌ی پایتون:} \lr{3.10}
    \item \textbf{نسخه‌ی پای‌تورچ:} \lr{2.13.0+cu118}
\end{itemize}

\subsection{تحلیل پیچیدگی و زمان آموزش}
در ادامه، خلاصه‌ای از وضعیت محاسباتی و عملکردی دو استراتژی اصلی که به‌طور کامل آموزش دیدند، آورده شده است. حجم فایل ذخیره‌شده‌ی مدل‌ها در هر دو حالت حدود \lr{43 MB} بوده است.

\begin{itemize}
    \item \textbf{مدل پایه \lr{(Frozen ResNet18)}:} \\
    پارامتر کل: \lr{11,180,103} | پارامتر قابل آموزش: \lr{3,591} | زمان کل آموزش (۵ دوره): \lr{10} دقیقه و \lr{14} ثانیه | میانگین زمان هر دوره: \lr{123} ثانیه | دقت نهایی روی آزمون: \lr{42\%} 
    
    \item \textbf{مدل ارتقایافته \lr{(Fine-Tuned ResNet18)}:} \\
    پارامتر کل: \lr{11,180,103} | پارامتر قابل آموزش: حدود \lr{8.4} میلیون | زمان کل آموزش (۳۰ دوره): \lr{63} دقیقه و \lr{12} ثانیه | میانگین زمان هر دوره: \lr{126} ثانیه | دقت نهایی روی آزمون: \lr{66\%} 
\end{itemize}
\newpage
\subsection{بحث و نتیجه‌گیری نهایی \lr{(Trade-off Analysis)}}
با بررسی داده‌های فوق، می‌توان به پرسش‌های کلیدی پروژه به این شکل پاسخ داد:
\begin{enumerate}
    \item \textbf{تأثیر تعداد پارامترها بر زمان:} با وجود اینکه در مدل پایه بیشتر لایه‌ها فریز بودند، زمان هر دوره (حدود \lr{123} ثانیه) تفاوت فاحشی با مدل \lr{Fine-Tuned} (حدود \lr{126} ثانیه) نداشت. دلیل این امر، وجود گلوگاه \lr{(Bottleneck)} در خواندن داده‌ها و انجام عملیات \lr{Data Augmentation} روی پردازنده‌ی مرکزی \lr{(CPU)} بود. با این حال، کاهش پارامترهای درگیر در مدل فریز شده، میزان مصرف حافظه‌ی گرافیکی \lr{(VRAM)} را به شدت کاهش داد.
    \item \textbf{یادگیری انتقالی در برابر آموزش از پایه \lr{(Scratch)}:} مدل دارای وزن‌های از پیش آموزش‌دیده، حتی در دوره‌ی اول به دقتی معادل \lr{53\%} رسید، در حالی که مدل طراحی شده از صفر \lr{(Scratch)} برای رسیدن به این دقت نیازمند ده‌ها دوره آموزش است. بنابراین یادگیری انتقالی بسیار کارآمدتر و سریع‌تر همگرا می‌شود.
    \item \textbf{توجیه عملی و انتخاب نهایی:} مدل فریز شده با حدود \lr{10} دقیقه آموزش به دقت \lr{42\%} رسید که برای کاربردهای واقعی و تجاری بسیار پایین است. در مقابل، اختصاص حدود یک ساعت زمان اضافی برای مدل \lr{Fine-Tuned} منجر به افزایش \lr{24} درصدی دقت (رسیدن به \lr{66\%}) شد. از آنجا که این افزایش دقت برای مسئله‌ی پیچیده‌ای مانند تشخیص احساسات بسیار حیاتی است، این هزینه‌ی محاسباتی اضافی کاملاً توجیه می‌گردد. 
    
    \textbf{نتیجه‌گیری:} مدل \lr{Fine-Tuned ResNet18} بهترین تعادل را بین زمان آموزش، پیچیدگی معماری و توانایی تعمیم‌پذیری ارائه می‌دهد و به عنوان مدل نهایی و برگزیده‌ی این پروژه معرفی می‌گردد.
\end{enumerate}